\SixthArticle{35}
But not only the product $\Pi\lambda.\Psi\lambda$, but also the function $\Psi\lambda$ itself, can be represented by a definite integral. With $k$ denoting a positive integer, it is clear that the value of the integral $\int\frac{x^\lambda-x^{\lambda+k}}{1-x}.\mathrm{d}x$ from $x=0$ to $x=1$ is
\[
=\frac{1}{\lambda+1}+\frac{1}{\lambda+2}+\frac{1}{\lambda+3}+\text{etc.}+\frac{1}{\lambda+k}
\]
Further, since the value of the integral $\int\left(\frac{1}{1-x}-\frac{kx^{k-1}}{1-x^k}\right)\mathrm{d}x$ is in general $=\mathrm{Const.}+\log\frac{1-x^k}{1-x}$, its value between the limits $x=0$ and $x=1$ will be $=\log k$; hence it is clear that the value of the integral $S=\int\left(\frac{1-x^\lambda+x^{\lambda+k}}{1-x}-\frac{kx^{k-1}}{1-x^k}\right)\mathrm{d}x$ between the same limits is
\[
=\log k-\frac{1}{\lambda+1}-\frac{1}{\lambda+2}-\frac{1}{\lambda+3}-\text{etc.}-\frac{1}{\lambda+k}
\]
an expression which we shall denote by $\Omega$. Let us split the integral $S$ into two parts:
\[
\int\left(\frac{1-x^\lambda}{1-x}\right)\mathrm{d}x+\int\left(\frac{x^{\lambda+k}}{1-x}-\frac{kx^{k-1}}{1-x^k}\right)\mathrm{d}x
\]
The first part, $\int\frac{1-x^\lambda}{1-x}.\mathrm{d}x$, becomes, on setting $x=y^k$,
\[
\int\frac{ky^{k-1}-ky^{\lambda k+k-1}}{1-y^k}\mathrm{d}y
\]
from which it follows of itself that its value from $x=0$ to $x=1$ equals the value of the integral
\[
\int\frac{kx^{k-1}-kx^{\lambda k+k-1}}{1-x^k}\mathrm{d}x
\]
between the same limits, since evidently the letter $y$ may be changed to $x$ under this restriction. Thus the integral $S$, between the same limits, becomes
\[
=\int\left(\frac{x^{\lambda+k}}{1-x}-\frac{kx^{\lambda k+k-1}}{1-x^k}\right)\mathrm{d}x
\]
This integral, however, on setting $x^k=z$, becomes
\SourceLeaf{168}
\[
\int\left(\frac{z^{\frac{\lambda+1}{k}}}{k(1-z)^{\frac{1}{k}}}-\frac{z^\lambda}{1-z}\right)\mathrm{d}z
\]
which, taken between the limits $z=0$ and $z=1$, therefore equals $\Omega$. But as $k$ grows to infinity, the limit of $\Omega$ is $\Psi\lambda$, the limit of $\frac{\lambda+1}{k}$ is $0$, while the limit of $k(1-z)^{\frac{1}{k}}$ is $\log\frac{1}{z}$ or $-\log z$. We therefore have
\[
\begin{aligned}
&[77]\qquad \Psi\lambda\\[4pt]
&=\int\left(\frac{1}{\log\frac{1}{z}}-\frac{z^\lambda}{1-z}\right)\mathrm{d}z
=\int\left(-\frac{1}{\log z}-\frac{z^\lambda}{1-z}\right)\mathrm{d}z
\end{aligned}
\]
from $z=0$ to $z=1$.

\SixthArticle{36}
The definite integrals by which the functions $\Pi\lambda$, $\Pi\lambda.\Psi\lambda$ were expressed above had to be restricted to values of $\lambda$ such that $\lambda+1$ is a positive quantity. This restriction followed from the derivation itself; indeed, it is easy to see that for other values of $\lambda$ those integrals always become infinite, even though the functions $\Pi\lambda$, $\Pi\lambda.\Psi\lambda$ may remain finite. For formula 77 to be true, the same condition, that $\lambda+1$ be a positive quantity, must certainly hold (otherwise the integral certainly becomes infinite, even though the function $\Psi\lambda$ remains finite); yet at first sight the derivation of the formula seems general and subject to no restriction. Closer attention, however, readily shows that this restriction is already present in the analysis by which the formula was obtained. For we tacitly assumed that the integral $\int\frac{1-x^\lambda}{1-x}\mathrm{d}x$, for which we substituted the equal integral $\int\frac{kx^{k-1}-kx^{\lambda k+k-1}}{1-x^k}\mathrm{d}x$, has a \emph{finite} value; this condition requires $\lambda+1$ to be a positive quantity. Our analysis does indeed show that these two integrals are always equal if the latter is taken from $x=0$ to $x=1-\omega$, the former from $x=0$ to $x=(1-\omega)^k$, however small the quantity $\omega$ may be, provided it is not $=0$. Nevertheless, when $\lambda+1$ is not a positive quantity, the two integrals, taken from $x=0$ to the \emph{same} endpoint $x=1-\omega$, do not approach equality at all; rather, their difference will then grow without bound as $\omega$ diminishes towards zero. This example shows how much circumspection is needed in dealing with infinite quantities, which in our judgment should be admitted into analytical reasoning only insofar as they can be reduced to the theory of limits.
\SourceLeaf{169}
\SixthArticle{37}
On setting $z=e^{-u}$ in formula 77, it is clear that it can also be written as follows:
\[
\begin{aligned}
&\Psi\lambda=-\int\left(\frac{e^{-u}}{u}-\frac{e^{-u\lambda-u}}{1-e^{-u}}\right)\mathrm{d}u,
\\[4pt]
&\quad\text{from }u=\infty\text{ to }u=0,\text{ i.e.}
\end{aligned}
\]
\[
\begin{aligned}
&\text{[78]}\qquad
\Psi\lambda=\int\left(\frac{e^{-u}}{u}-\frac{e^{-\lambda u}}{e^u-1}\right)\mathrm{d}u,
\\[4pt]
&\quad\text{from }u=0\text{ to }u=\infty.
\end{aligned}
\]
(Similarly, the value of $\Pi\lambda$ given in Article 28 becomes, on setting $e^{-y}=\upsilon$, the following:
\[
\Pi\lambda=\int\left(\log\frac{1}{\upsilon}\right)^\lambda\mathrm{d}\upsilon,
\quad\text{from }\upsilon=0\text{ to }\upsilon=1)
\]
Further, it follows from formula 77 that
\[
\text{[79]}\qquad
\Psi\lambda-\Psi\mu=\int\frac{z^\mu-z^\lambda}{1-z}\mathrm{d}z,
\quad\text{from }z=0\text{ to }z=1
\]
where, besides $\lambda+1$, $\mu+1$ must also be a positive quantity.

On setting $z=u^\alpha$ in the same formula 77, with $\alpha$ denoting a positive quantity, the result is
\[
\begin{aligned}
&\Psi\lambda=\int\left(-\frac{u^{\alpha-1}}{\log u}
-\frac{\alpha u^{\alpha\lambda+\alpha-1}}{1-u^\alpha}\right)\mathrm{d}u,
\\[4pt]
&\quad\text{from }u=0\text{ to }u=1
\end{aligned}
\]
and since, similarly, for a positive value of $\beta$ it is possible to set
\[
\Psi\lambda=\int\left(-\frac{u^{\beta-1}}{\log u}
-\frac{\beta u^{\beta\lambda+\beta-1}}{1-u^\beta}\right)\mathrm{d}u
\]
it is clear that
\[
0=\int\left(\frac{u^{\alpha-1}-u^{\beta-1}}{\log u}
+\frac{\alpha u^{\alpha\lambda+\alpha-1}}{1-u^\alpha}
-\frac{\beta u^{\beta\lambda+\beta-1}}{1-u^\beta}\right)\mathrm{d}u
\]
or
\[
\int\frac{u^{\alpha-1}-u^{\beta-1}}{\log u}\mathrm{d}u
=\int\left(\frac{\beta u^{\beta\lambda+\beta-1}}{1-u^\beta}
-\frac{\alpha u^{\alpha\lambda+\alpha-1}}{1-u^\alpha}\right)\mathrm{d}u
\]
with the integrals always taken from $u=0$ to $u=1$. But on setting $\lambda=0$, the latter integral can be given \emph{indefinitely}; namely, it is $=\log\frac{1-u^\alpha}{1-u^\beta}$, if it is to vanish for $u=0$. Thus, since for $u=1$ it is necessary to set $\frac{1-u^\alpha}{1-u^\beta}=\frac{\alpha}{\beta}$, the integral will be $\log\frac{\alpha}{\beta}=\int\frac{u^{\alpha-1}-u^{\beta-1}}{\log u}\mathrm{d}u$, from $u=0$ to $u=1$, a theorem which the distinguished \textsc{Euler} once obtained by other methods.
\par\medskip\noindent\hfill\rule{.25\linewidth}{.4pt}\hfill\par
